\chapter{Photovoltaic System}
\label{ch:pv}

\section{System Specifications}

The reference PV system is a 4~kWp rooftop installation with the parameters in Table~\ref{tab:pv_specs}.

\begin{table}[H]
\centering
\caption{PV system parameters.}
\label{tab:pv_specs}
\begin{tabular}{lll}
\toprule
\textbf{Parameter} & \textbf{Value} & \textbf{Unit} \\
\midrule
Peak power & 4,000 & Wp \\
Panel tilt & 30 & degrees \\
Azimuth & 0 (due south) & degrees \\
Temperature coefficient & $-0.004$ & /°C \\
Inverter efficiency & 96 & \% \\
STC irradiance & 1,000 & W/m$^2$ \\
STC temperature & 25 & °C \\
\bottomrule
\end{tabular}
\end{table}

\section{PV Power Model}

The PV power output is modeled as:

\begin{equation}
    P_{\text{PV}}(t) = P_{\text{peak}} \cdot \frac{G(t)}{G_{\text{STC}}} \cdot \left(1 + \gamma \cdot (T_{\text{cell}}(t) - T_{\text{STC}})\right) \cdot \eta_{\text{inv}}
    \label{eq:pv_power}
\end{equation}

where:
\begin{itemize}[nosep]
    \item $P_{\text{peak}}$ = rated PV capacity (4,000~W)
    \item $G(t)$ = global horizontal irradiance at time $t$ (W/m$^2$)
    \item $G_{\text{STC}}$ = 1,000~W/m$^2$
    \item $\gamma$ = temperature coefficient ($-0.004$~/°C)
    \item $T_{\text{cell}}(t)$ = cell temperature (°C)
    \item $T_{\text{STC}}$ = 25~°C
    \item $\eta_{\text{inv}}$ = inverter efficiency (0.96)
\end{itemize}

The cell temperature is estimated using the simplified NOCT model:

\begin{equation}
    T_{\text{cell}}(t) = T_{\text{ambient}}(t) + 0.03 \cdot G(t)
    \label{eq:tcell}
\end{equation}

\section{Clear-Sky Irradiance Model}

For the simulation, clear-sky irradiance is computed using solar geometry (declination, hour angle, solar elevation) and a simplified atmospheric transmittance model. Cloud cover is applied as a multiplicative factor:

\begin{equation}
    G(t) = G_{\text{clear}}(t) \cdot (1 - 0.75 \cdot C_{\text{cloud}})
\end{equation}

where $C_{\text{cloud}} \in [0, 1]$ represents fractional cloud cover.
